\section{libDAI: exact factor meaning before approximate computation}
\label{sec:exec24-libdai-reuse}

A bank may need probabilities from an already specified causal model before
it compares interventions. A discrete inference engine can compute those
probabilities, but it cannot supply missing causal assumptions merely by
returning normalized beliefs. libDAI is a useful reference for that boundary:
its factor representation, exact enumeration, junction-tree computation and
belief propagation expose different mathematical and operational contracts.
The earlier combined C++ comparison remains intact; this review separates
libDAI's concrete reuse decision from OpenGM and examines the selected
implementation rather than treating a feature list as qualification
\citep{exec13_cxx_inference,exec24_libdai_source}.

\subsection{A historical paper and a later source revision}
Mooij's five-page JMLR paper describes version 0.2.7, released under GPL v2 or
later. Its factor-graph formulation, available inference tasks and interface
scope are primary technical evidence. Its comparison table is historical,
and its performance statements are not measurements of the proposed banking
worker. The genuine publisher PDF is retained and inspected, including its
distinction among marginals, partition functions and maximum-probability
configurations. The paper says that not every method implements every task;
it does not promise a general causal-discovery or identification service
\citep{exec24_libdai_paper}.

The selected source is the \texttt{dbtsai/libDAI} mirror at its 16 May 2012
revision. Its README declares version 0.3.0, dated 12 July 2011, and its
license grants redistribution under two-clause BSD conditions. Those later
terms must not be projected backwards onto the paper's version. Nineteen
selected files authenticate against the complete pinned tree, including
factor and probability primitives, indexing, factor-graph operations, exact
inference, belief propagation, junction-tree construction and build settings.
The mirror is not authenticated as an official current distribution. The
official-source attempt was rejected by the collector's fixed host policy,
which is not evidence of an unavailable server
\citep{exec24_libdai_source,exec24_libdai_choice}.

The Linux template links GMP and GMPXX and names Boost libraries. Its SWIG
example points to Python 2.7 headers. Those are concrete porting questions,
not proof that a modern Python package is unusable or already supported.
The platform-independent configuration enables debug checks, while the
makefile permits another configuration. A qualified artifact must record
the actual build options and complete dependency notices. The inspected BSD
notice is evidence about the selected source, not complete clearance of a
future binary and its bindings. No downloaded build script, source or native
library has been executed for this review \citep{exec24_libdai_source}.

The preceding bridge and discovery examples concern the interpretation of a
graph inferred or supplied by an analyst. Even after a causal model is
accepted for a particular calculation, its probability representation can
change at the next boundary. Factor ordering, normalization and the operation
used to encode an intervention therefore require their own comparison.

\subsection{Factor meaning is a typed conversion contract}
Let the finite variables be $X_V=(X_i:i\in V)$ and let each nonnegative factor
$f_I$ depend on variables $X_{N_I}$. The represented probability law is
\[
 p(x_V)=\frac{\prod_I f_I(x_{N_I})}{Z},\qquad
 Z=\sum_{x_V}\prod_I f_I(x_{N_I}).
\]
This requires finite nonnegative entries and $0<Z<\infty$. A Bayesian network
can be represented using one conditional-probability factor per variable.
That conversion preserves a probability law only if category meaning,
variable identity, dimensions and entry ordering are also preserved
\citep{exec24_libdai_paper,exec24_libdai_source}.

The source orders variables by ascending numeric label and uses the first
variable as the fastest-changing index. With binary $X$ followed by
three-state $Y$, the index is $x+2y$. Assigning weights $(1,2,3,4,5,6)$ in
that order yields $p(X=1)=4/7$. Reading the same vector as a row-major table
with index $3x+y$ instead yields $5/7$. Both outputs normalize and can look
plausible; a probability-simplex check cannot detect this error. One factor
constructor accepts a variable set and copies entries, whereas the ordered
variable-vector constructor explicitly permutes entries. A proposed adapter
must select the correct contract and retain category labels and the
permutation as auditable input \citep{exec24_libdai_source,exec24_libdai_choice}.

The variable comparison operators use labels for identity; matching labels
with different state counts trigger checks only under the selected debug
conditional. The adapter should always reject conflicting domains before
construction, regardless of the native build. The linear-index routine also
defaults unspecified variables to state zero. A missing business observation
must therefore be represented as missing evidence, not passed as an omitted
state that silently becomes the first category. The marginal routine
intersects the requested variable set with the factor's domain; the caller
must verify that every requested variable was actually represented. These
are distinct identity, missingness and query-completeness tests
\citep{exec24_libdai_source}.

\subsection{Clamping evidence does not perform a causal intervention}
The inspected factor-graph clamp multiplies neighboring factors by a
single-state indicator. This conditions the original distribution on an
observation. In a causally justified Bayesian network, an intervention
instead replaces the factor governing the action by the specified assignment
mechanism, while retaining the other mechanisms. The factor graph alone
does not say which table was an action's causal mechanism. That provenance
must survive compilation \citep{exec24_libdai_source,exec24_libdai_choice}.

An exact authored example makes the distinction observable. Let $U$ be fair
binary, let $P(A=U\mid U)=4/5$, and let $Y=U$. Then
$P(Y=1\mid A=1)=4/5$, while
$P(Y=1\mid\operatorname{do}(A=1))=1/2$ and the intervention contrast is zero.
The conditioning calculation retains $p(A\mid U)$; the intervention
calculation replaces it. This example supplies the causal model explicitly.
It does not infer that model from the observed distribution
\citep{exec24_libdai_choice}.

Nor is making a cavity a generic intervention implementation: the selected
method replaces every factor touching a variable with ones, which can erase
outcome mechanisms as well as the assignment mechanism. Backup and restore
operations can retain earlier factors, but the worker must verify recovery
under failed or cancelled requests. In the selected clamp, the initial
guard permits an index equal to the number of states, whereas valid indices
are strictly smaller. The subsequent setter has a debug-only bounds check.
This is a conditional source boundary requiring a native fixture, not an
observed memory fault. The proposed API rejects every index outside
$0\leq x<r_i$ before calling the library
\citep{exec24_libdai_source,exec24_libdai_choice}.

\subsection{Exact inference has numerical and resource conditions}
The selected exact enumerator multiplies factors into a full joint table,
computes its sum and normalizes marginals. Its algorithmic exactness is
conditional on the supplied finite model; it does not imply exact arithmetic.
For two singleton factors with entries $(10^{-200},2\cdot10^{-200})$, the
mathematically normalized product is $(1/5,4/5)$, but both binary64 products
underflow to zero in the independent witness. A log-scale calculation with
maximum subtraction recovers the intended normalized ratio. An exactly zero
normalizer instead denotes an impossible or invalid query and must be
refused. Scaling factors and tracking their scales can preserve normalized
beliefs, but the scales must also be restored when reporting $\log Z$.
The inspected probability primitive explicitly refuses a zero norm;
finite, nonnegative inputs and finite outputs require additional checks
\citep{exec24_libdai_source,exec24_libdai_choice}.

Resource use is another part of the result. A full joint table on thirty
binary variables has $2^{30}$ entries: eight-byte entries alone require
8 GiB, before containers and intermediate products. Junction-tree inference
can exploit structure, but its clique tables may still be large. The
selected constructor chooses an elimination heuristic and applies a memory
estimate with a documented heuristic multiplier. An estimate is not an
operating-system resource limit. A proposed worker needs its own memory,
time, cancellation and output limits, and must report when a requested
exact calculation was refused or replaced by an explicitly approved
approximation \citep{exec24_libdai_source,exec24_libdai_choice}.

\subsection{A zero convergence residual can coexist with a wrong marginal}
Belief propagation iterates messages, with configured update order, damping,
log-domain arithmetic, tolerance and iteration/time bounds. The selected
implementation returns the largest change in beliefs; a nonconvergence
warning is conditional on verbosity. It explicitly omits NaN checks in the
iteration loop. A worker must inspect residuals, iteration counts, finite
values and termination conditions independently of textual logs. Even a
finite residual below tolerance does not bound marginal error
\citep{exec24_libdai_source}.

Consider an authored triangle of binary variables with edge factor three
when neighboring states agree and one when they disagree. Uniform messages
remain uniform after a belief-propagation update, giving zero residual.
The resulting edge belief assigns agreement probability $3/4$. Exact
enumeration gives two all-equal states of weight 27 each and six other
states of weight three each, hence $Z=72$. For a fixed edge, the agreeing
weight is 60, so its true agreement probability is $5/6$. The error is
$1/12$ despite the fixed point. Single-variable marginals are one half in
both calculations, so inspecting only those marginals would miss the
discrepancy \citep{exec24_libdai_choice}.

This is a mathematical witness for the proposed diagnostic, not a native
libDAI benchmark. It distinguishes iterative stability, approximation error
and statistical uncertainty. Under a hypothetical decision coefficient of
120,000 USD, an error of $1/12$ translates into 10,000 USD of value error.
A small residual cannot justify that monetary exposure. Public qualification
should compare exact and approximate results on declared finite regimes,
measure query-specific errors and refuse unsupported accuracy claims.
Sampling uncertainty and uncertainty about the causal model remain separate
\citep{exec24_libdai_choice}.

\subsection{A reference choice with a concrete integration route}
The current decision is to retain libDAI as a historical numerical reference
and conditional candidate for an isolated C++ inference worker. Direct
adoption of a general causal platform is rejected because the selected
surface does not supply identification. A port is deferred until a measured
need justifies a new correctness and maintenance burden. Reuse through a
rights-reviewed worker can be preferable to porting when exact artifact
identity, factor semantics and error contracts are demonstrably preserved.
Another current inference implementation is an alternative comparator,
but it must face the same typed queries and public oracles; this review does
not confer approval on an uninspected substitute \citep{exec24_libdai_choice}.

The known route is to preregister finite fixtures for variable permutations,
conflicting cardinalities, unspecified states and absent query variables;
conditioning versus intervention and restoration after faults; zero mass,
nonfinite entries and underflow; exact enumeration versus junction trees;
tree and loopy-message fixed points, truncation and resource refusal.
The numerical owner supplies exact oracles and error budgets; the causal
owner supplies mechanism provenance and identifying assumptions; the
adapter owner supplies types and refusal behavior; the platform/rights owner
supplies artifact and dependency qualification. Every submitted request,
including conversion and runtime failures, remains in the cost denominator.
A purely illustrative 12,000 USD over 400 submitted, 300 completed and
200 qualified requests produces 30, 40 and 60 USD respectively.
Official provenance, applicable rights, execution, independent review, current support and actual costs remain unresolved \citep{exec24_libdai_choice}.
